\chapter{Database Security and Administration}
\label{chap:security}

\section{User Roles and Permissions}

Three database user roles are defined following the principle of least privilege:

\begin{table}[H]
    \centering
    \caption{Database User Roles and Permissions}
    \begin{tabular}{llp{7cm}}
        \toprule
        \textbf{Role} & \textbf{Username} & \textbf{Permissions} \\
        \midrule
        Administrator & \texttt{db\_admin} & Full privileges: CREATE, DROP, ALTER, INSERT, UPDATE, DELETE, SELECT, GRANT \\
        Analyst & \texttt{db\_analyst} & SELECT on all tables, EXECUTE on stored procedures and UDFs \\
        App User & \texttt{db\_app} & SELECT, INSERT on Transactions and BehaviorSignals; EXECUTE procedures \\
        \bottomrule
    \end{tabular}
\end{table}

\begin{lstlisting}[style=sqlstyle, caption={Role Configuration}]
-- Create analyst role (read-only access for reporting)
CREATE USER 'db_analyst'@'localhost' IDENTIFIED BY 'analyst_pass';
GRANT SELECT ON fomo_project.* TO 'db_analyst'@'localhost';
GRANT EXECUTE ON PROCEDURE fomo_project.sp_refresh_behavior_signals 
    TO 'db_analyst'@'localhost';

-- Create app user (limited write access)
CREATE USER 'db_app'@'localhost' IDENTIFIED BY 'app_pass';
GRANT SELECT, INSERT ON fomo_project.Transactions TO 'db_app'@'localhost';
GRANT SELECT ON fomo_project.vw_warning_dashboard TO 'db_app'@'localhost';
GRANT EXECUTE ON PROCEDURE fomo_project.sp_simulate_new_transactions 
    TO 'db_app'@'localhost';

FLUSH PRIVILEGES;
\end{lstlisting}

\section{Data Security Considerations}

Although this project uses a synthetic dataset, the following security measures are proposed for a production deployment:

\begin{itemize}
    \item \textbf{Encryption at rest:} Enable MySQL InnoDB tablespace encryption for tables containing investor profiles.
    \item \textbf{Encryption in transit:} Enforce TLS/SSL for all client-server connections.
    \item \textbf{Password policies:} Enforce strong password requirements using MySQL \texttt{validate\_password} plugin.
    \item \textbf{Connection string security:} Store credentials in environment variables via \texttt{.env} files (implemented in \texttt{config.py}), never hardcoded.
    \item \textbf{Audit logging:} Enable MySQL General Query Log for administrative actions.
\end{itemize}

\section{Backup and Recovery}

\begin{lstlisting}[language=bash, caption={Backup Procedure using mysqldump}]
# Full database backup
mysqldump -u db_admin -p fomo_project \
    --single-transaction \
    --routines \
    --triggers \
    > backup_fomo_$(date +%Y%m%d).sql

# Restore from backup
mysql -u db_admin -p fomo_project < backup_fomo_20260101.sql
\end{lstlisting}

A recommended backup schedule:
\begin{itemize}
    \item Full backup: weekly
    \item Incremental backup (binary log): daily
    \item Retention period: 30 days
\end{itemize}

\section{Performance Optimization}

Beyond indexing (covered in Chapter~\ref{chap:advanced}), the following optimizations are applied:

\begin{itemize}
    \item \textbf{InnoDB buffer pool:} Sized to accommodate the full Portfolio\_History table in memory for fast range scans.
    \item \textbf{Query caching via views:} Frequently accessed join queries are encapsulated in views to simplify the query planner's work.
    \item \textbf{Chunked inserts:} Simulation data is inserted in batches of 1,000 rows using SQLAlchemy's \texttt{chunksize} parameter.
    \item \textbf{EXPLAIN analysis:} All major queries were verified with \texttt{EXPLAIN} to confirm index usage and avoid full table scans.
\end{itemize}
